% Generated draft: requires mathematical and source review.
\documentclass[11pt]{article}
\usepackage[margin=25mm]{geometry}
\usepackage{fontspec}
\setmainfont{Latin Modern Roman}
\usepackage{amsmath,amssymb,mathtools,unicode-math}
\setmathfont{Latin Modern Math}
\usepackage{enumitem,needspace,xurl,hyperref}
\setlength{\parindent}{0pt}
\setlength{\parskip}{5pt}
\setlength{\emergencystretch}{3em}
\newcommand{\N}{\mathbb{N}}
\newcommand{\Z}{\mathbb{Z}}
\newcommand{\Q}{\mathbb{Q}}
\newcommand{\R}{\mathbb{R}}
\newcommand{\C}{\mathbb{C}}
\newcommand{\F}{\mathbb{F}}
\newcommand{\A}{\mathbb{A}}
\newcommand{\PP}{\mathbb{P}}
\newcommand{\E}{\mathbb{E}}
\newcommand{\eps}{\varepsilon}
\newcommand{\dd}{\,\mathrm{d}}
\begin{document}
\section*{UnsolvedMath ID: 30000625}
\textbf{Source problem code:} catalogue entry
\subsection{Definitions and mathematical statement}
Let \(F\) be a field and \(A\) a finite-dimensional central simple \(F\)-algebra: its center is \(F\) and it has no nonzero proper two-sided ideals. Its degree is \(\deg A=\sqrt{\dim_F A}\). The source attributes to Albert the conjecture that every such algebra of prime degree \(p\) is cyclic, with no restriction on \(F\) or its characteristic stated.

For a division algebra \(D\) of degree \(p\), the standard cyclicity assertion means that it has a maximal subfield \(L\subseteq D\) which is a cyclic Galois extension of \(F\) of degree \(p\). Equivalently, the requested presentation is a cyclic crossed product
\[
 D\cong(L/F,\sigma,a)=\bigoplus_{i=0}^{p-1}Lu^i,
 \qquad u\ell=\sigma(\ell)u,\quad u^p=a,
\]
where \(\operatorname{Gal}(L/F)=\langle\sigma\rangle\) and \(a\in F^\times\). The conjectural nonsplit content is therefore
\[
 \forall F\ \forall p\text{ prime}\ \forall D\text{ central division over }F,
 \qquad \deg D=p\ \Longrightarrow\ D\text{ cyclic}.
\]
Cyclic refers to the Galois group of a maximal splitting subfield, not to the additive or multiplicative group of the algebra. No primitive-root-of-unity hypothesis or Kummer presentation is inserted, so the formulation also leaves characteristic \(p\) within the source's scope.

A prime-degree central simple algebra is either division or split, the latter isomorphic to \(\operatorname{Mat}_p(F)\). The word cyclic for a split algebra needs a convention: one often counts the split case automatically, or permits a cyclic \(\acute{e}\)tale algebra such as \(F^p\) in a crossed-product formulation. If one insists instead on an actual degree-\(p\) cyclic field extension for every split algebra, a field without such an extension creates an irrelevant additional obstruction. The source does not fix that convention, so the nontrivial division formulation is distinguished explicitly without discarding the original central-simple wording.

The literature triage checked 2026-08-21 reports no authoritative complete resolution of the extracted conjecture located in its search. It gives no prime-degree special-case theorem or field restrictions. The open label and editorial verification note remain dated supplied metadata, not independently checked current conclusions.
\subsection{Short English statement}
Albert's conjecture asks whether every prime-degree central division algebra has a cyclic Galois maximal subfield, with split central simple algebras handled by the intended cyclicity convention. No field-characteristic or roots-of-unity restriction is added. The split-case terminology is made explicit rather than imposing existence of a field extension where none may exist.
\subsection{Sources}
\begin{itemize}
\item Local UnsolvedMath registry, permanent ID 30000625, source code catalogue entry, ``Albert's Cyclicity Conjecture.'' The original wording attributes the prime-degree central-simple assertion to Albert.
\item \emph{Infinite-Dimensional Lie Theory (2007)}, Oberwolfach Reports, pp.\ 3259--3320, DOI 10.4171/owr/2006/55: \url{https://doi.org/10.4171/owr/2006/55}. This is also the sole supplied literature reference. The citation uses 2007, while the DOI and proposed-year metadata use 2006.
\item Literature triage checked 2026-08-21; record updated 2026-09-16. No external source has been visited or independently checked here.
\end{itemize}
\end{document}
